%% ma: L12-B08-0165
%% lop: 12
%% chuong: 2
%% bai: 8
%% dang: 12.8.52
%% loai: TN4
%% muc: VD
%% dap_an: "D"
%% nguon: "CĐ2.4.pdf giáo viên gửi 27/10/2026 (M07, Lê Minh Khôi – Nguyễn Anh Khôi, Chương II Chủ đề 4 Các bài toán về tâm tỉ cự), Câu 15"
%% kien_thuc: "Bài 8 (biểu thức tọa độ của tích vô hướng), lớp 10 (giá trị lớn nhất của tam thức bậc hai)"
%% trang_thai: chinh-thuc
%% ghi_chu: "Giáo viên đã duyệt 27/10/2026"
\begin{ex}
Trong không gian $Oxyz$, cho ba điểm $A(4;1;5)$, $B(3;0;1)$, $C(-1;2;0)$ và điểm $K(a;b;c)$ thỏa mãn $\overrightarrow{KA}\cdot\overrightarrow{KB}+2\overrightarrow{KB}\cdot\overrightarrow{KC}-5\overrightarrow{KC}\cdot\overrightarrow{KA}$ đạt giá trị lớn nhất. Tính giá trị của biểu thức $P=a-2b+4c$.
\choice{$P=23$}{$P=31$}{$P=11$}{\True $P=13$}
\loigiai{Biểu thức bằng $-2a^2+4a-2b^2+10b-2c^2+17c+21$, lớn nhất khi $a=1$, $b=\dfrac52$, $c=\dfrac{17}4$.

$P=1-5+17=13$.}
\end{ex}
